% Сборка из корня репозитория:
% lualatex -halt-on-error -interaction=nonstopmode -output-directory=SVD SVD/SVD_solver.tex
% Команду следует выполнить дважды для оглавления и перекрёстных ссылок.
\documentclass[11pt,a4paper]{article}
\usepackage{fontspec}
\setmainfont{Noto Serif}
\setsansfont{Noto Sans}
\setmonofont[Scale=0.82]{DejaVu Sans Mono}
\usepackage[russian]{babel}
\usepackage[margin=24mm]{geometry}
\usepackage{amsmath,amssymb,amsthm,mathtools}
\usepackage{booktabs,tabularx,longtable,array}
\usepackage{enumitem,microtype,xurl}
\usepackage[unicode,hidelinks]{hyperref}
\usepackage{bookmark}
\hypersetup{pdftitle={SVD-решатель для многоиндексной процедуры средних производных},
  pdfauthor={},pdfsubject={Теория, ограничения и направления развития}}
\setlist{itemsep=3pt,topsep=5pt}
\setlength{\emergencystretch}{2em}
\setlength{\parskip}{3pt}
\numberwithin{equation}{section}
\newtheorem{theorem}{Теорема}[section]
\newtheorem{lemma}[theorem]{Лемма}
\newtheorem{proposition}[theorem]{Утверждение}
\theoremstyle{remark}
\newtheorem*{remark}{Замечание}
\newcommand{\R}{\mathbb R}
\newcommand{\Id}{\mathrm{Id}}
\newcommand{\T}{^{\mathsf T}}
\newcommand{\F}{_{\mathrm F}}
\newcommand{\ip}[2]{\langle #1,#2\rangle_{\mathrm F}}
\DeclareMathOperator{\rank}{rank}
\DeclareMathOperator{\tr}{tr}
\DeclareMathOperator{\diag}{diag}
\DeclareMathOperator{\vecop}{vec}
\DeclareMathOperator*{\argmin}{arg\,min}
\DeclareMathOperator*{\argmax}{arg\,max}
\newcommand{\src}[1]{\nolinkurl{#1}}
\newcommand{\proposal}{\par\noindent\textbf{Предложение для исследования.}\ }

\title{SVD-решатель для многоиндексной\allowbreak{} процедуры средних производных\\[5pt]
  \large Математическая постановка, ограничения\allowbreak{} и три направления развития}
\author{}
\date{29 сентября 2026 года}

\begin{document}
\maketitle
\begin{abstract}
Рассматривается внутренний шаг многоиндексной процедуры средних производных
(ADP), в котором локальные коэффициенты зафиксированы, а глобальная матрица
оценивается по взвешенному квадратичному функционалу с регуляризацией.
Выведены жадное построение компонент ранга один, совместная переоценка их
масштабов, точное решение в изотропной модели и проверяемые оценки ошибки.
Показано, почему усечение ранга самой матрицы может противоречить восстановлению
полного пространства эффективного снижения размерности (EDR).
Развитие метода организовано по трём направлениям: сочетание SVD со вторым
решателем, совместная обработка блоков на GPU, улучшение начального приближения и
выбора направлений. Для каждого направления приведены уравнения, условия
применимости и способ проверки. Новые предложения отделены от доказанных
следствий постановки и от результатов имеющейся реализации.
\end{abstract}
\tableofcontents
\newpage

\section{Задача и единая система обозначений}\label{sec:setup}

\subsection{Место решателя в ADP}
Пусть наблюдения $(x_i,y_i)$, $i=1,\ldots,n$, подчиняются модели
\[
  y_i=f_0(P_\circ x_i)+\varepsilon_i,
  \qquad P_\circ\in\R^{m\times d},\qquad P_\circ P_\circ\T=\Id_m.
\]
Искомый объект --- пространство строк $P_\circ$, а не конкретный выбор
базиса в нём. Его ортогональный проектор равен
$\Pi_\circ=P_\circ\T P_\circ$. Для восстановления пространства нужны
предположения о регрессионной функции, распределении данных и возбуждении
всех $m$ направлений. Одна лишь запись многоиндексной модели их не обеспечивает.

Во внешнем цикле ADP выбираются центры локализации, веса и пробные
направления, вычисляются локальные статистики и оценивается новый глобальный
базис. Здесь подробно изучается \emph{одна внутренняя задача} этого цикла.
Её исходные величины фиксируются на время решения. Изменение весов,
коэффициентов или направлений начинает другую внутреннюю задачу.

\begin{center}
\begin{tabularx}{\linewidth}{@{}l l X@{}}
\toprule
Обозначение & Размер & Смысл\\
\midrule
$d$, $m$ & $m\le d$ & Размерность признаков и искомого пространства\\
$J$, $p$ & целые числа & Число центров и число пробных направлений на центр\\
$P$ & $m\times d$ & Текущий базис, $PP\T=\Id_m$\\
$B$ & $m\times d$ & Матрица, оптимизируемая во внутренней задаче\\
$g_j$ & $m$ & Зафиксированные локальные коэффициенты, столбец\\
$I_j$ & $p$ & Вектор локальных откликов\\
$U_j$ & $p\times d$ & Локальный линейный оператор\\
$N_j$ & скаляр & Локальная масса, $N_j\ge0$\\
$\lambda$ & скаляр & Коэффициент регуляризации\\
$r$, $q$ & целые числа & Ограничение ранга и размер блока направлений\\
\bottomrule
\end{tabularx}
\end{center}

Буква $P$ обозначает матрицу базиса; число пробных направлений обозначено
строчной $p$, чтобы исключить неоднозначность исходных текстов.
В реализации все $U_j$ собраны в массив размера $(J,p,d)$.
Статистики $I_j,U_j$ должны иметь одну и ту же согласованную нормировку.
В нормированной версии ADP они уже разделены на локальную массу, а $N_j$
входит во внешний вес функционала. Повторное умножение статистик на $N_j$
создаёт другую задачу.

Теория допускает $N_j=0$: такой центр не влияет на сумму. Действующая
проверка входов решателя требует строго положительных конечных масс.
Эти два условия следует различать. Если не оговорено иное, далее
$\lambda>0$; случай $\lambda=0$ обсуждается отдельно.

\subsection{Три различные постановки}\label{sec:variants}
Общий квадратичный функционал имеет вид
\begin{equation}\label{eq:F}
 F(B)=\sum_{j=1}^J N_j\|I_j-U_jB\T g_j\|_2^2
       +\lambda\|B-P\|\F^2.
\end{equation}
Рассмотрим три допустимых множества:
\begin{align}
 &\min_{\rank(B)\le r} F(B),\quad r<m,
 &&\text{усечение ранга матрицы};\label{eq:matrix}\\
 &\min_{\rank(\Delta)\le r} F(P+\Delta),
 &&\text{малоранговая поправка};\label{eq:correction}\\
 &\min_{B\in\R^{m\times d}} F(B),
 &&\text{полный условный шаг}.\label{eq:full}
\end{align}
Первая постановка --- основная задача исходных \src{SVD.tex} и
\src{SVD_form.tex}. Вторая также представлена отдельным режимом действующего
решателя \src{ADP/solver/SVD.py}; она не приписывается исходным диалогам.
В ней малый ранг относится к изменению базиса,
поэтому $P+\Delta$ может иметь полный ранг $m$. В третьей ограничение ранга
снято. Различие допустимых множеств существенно даже при одинаковом
выражении функционала.

Обозначение SVD означает сингулярное разложение. Для общего набора $U_j$
решатель использует SVD как часть вычислений; сама задача \eqref{eq:matrix}
не сводится к обычному усечению сингулярных значений произвольной матрицы.
Кроме того, функционал \eqref{eq:F} нельзя без отдельной проверки отождествлять
со всеми вариантами HPAO, встречающимися в коде и рукописях.

\section{Квадратичная операторная форма}\label{sec:operator}

\subsection{Прямой и сопряжённый операторы}
На пространстве $\R^{m\times d}$ используем скалярное произведение
$\ip{X}{Y}=\tr(X\T Y)$. Введём
\begin{equation}\label{eq:A}
 (\mathcal A X)_j=\sqrt{N_j}\,U_jX\T g_j,
 \qquad b_j=\sqrt{N_j}\,I_j.
\end{equation}
Пространство значений --- набор $J$ векторов длины $p$ с обычным
суммарным скалярным произведением. Прямой расчёт даёт
\begin{equation}\label{eq:adjoint}
 \mathcal A^* z=\sum_j\sqrt{N_j}\,g_j(U_j\T z_j)\T.
\end{equation}
Действительно,
\[
 \sum_j z_j\T(\mathcal A X)_j
 =\sum_j\sqrt{N_j}\,g_j\T XU_j\T z_j
 =\ip{X}{\mathcal A^*z}.
\]
Эта тождественность должна проверяться численно при любой новой реализации,
включая перенос на GPU и разбиение центров на пакеты.

Положим
\begin{equation}\label{eq:LC}
 \mathcal L=\mathcal A^*\mathcal A+\lambda\Id,
 \qquad C=\mathcal A^*b+\lambda P.
\end{equation}
Если $H_j=U_j\T U_j$ и $c_j=U_j\T I_j$, то
\begin{align}
 \mathcal L(X)&=\sum_jN_j(g_jg_j\T)XH_j+\lambda X,\label{eq:L}\\
 C&=\sum_jN_jg_jc_j\T+\lambda P.\label{eq:C}
\end{align}
Матрицы $H_j$ нужны для вывода, но обычно не нужны для вычислений.
Применение $\mathcal L$ выполняется через $U_j$ и $U_j\T$.

\subsection{Положительная определённость и полный минимум}
\begin{lemma}\label{lem:spd}
Оператор $\mathcal L$ самосопряжён. При $\lambda>0$ он положительно
определён и удовлетворяет
\begin{equation}\label{eq:bounds}
 \lambda\|X\|\F^2\le\ip{X}{\mathcal L(X)}
 \le\beta\|X\|\F^2,
 \qquad
 \beta=\lambda+\sum_jN_j\|g_j\|_2^2\|U_j\|_2^2.
\end{equation}
\end{lemma}
\begin{proof}
Самосопряжённость следует из $\mathcal A^*\mathcal A$.
Также
\[
 \ip{X}{\mathcal L(X)}
 =\sum_jN_j\|U_jX\T g_j\|_2^2+\lambda\|X\|\F^2.
\]
Нижняя оценка получается удалением неотрицательной суммы, верхняя ---
из $\|U_jX\T g_j\|_2\le\|U_j\|_2\|X\|\F\|g_j\|_2$.
\end{proof}

Раскрывая квадраты, получаем
\begin{equation}\label{eq:quadratic}
 F(B)=C_0+\ip{B}{\mathcal L(B)}-2\ip{C}{B},
 \quad C_0=\sum_jN_j\|I_j\|_2^2+\lambda m.
\end{equation}
Таким образом,
\begin{equation}\label{eq:Q}
 \nabla F(B)=2(\mathcal L(B)-C),
 \qquad Q(B)=C-\mathcal L(B)=-\tfrac12\nabla F(B).
\end{equation}
$Q$ будем называть невязкой нормального уравнения. Она имеет размер
$m\times d$ и отличается от невязок данных
$e_j=I_j-U_jB\T g_j$.

Полный минимум $B_\star$ единственен и решает
\begin{equation}\label{eq:normal}
 \mathcal L(B_\star)=C.
\end{equation}
Для любого $B$ справедливо точное тождество
\begin{equation}\label{eq:energy}
 F(B)-F(B_\star)=\ip{B-B_\star}{\mathcal L(B-B_\star)}.
\end{equation}
Следовательно, задача \eqref{eq:matrix} является поиском лучшей матрицы
ранга не выше $r$ в норме, задаваемой $\mathcal L$.
В общем случае эта норма не является нормой Фробениуса.

При векторизации \emph{по столбцам} матрица оператора имеет вид
\begin{equation}\label{eq:kron}
 K=\sum_jN_j\bigl(H_j\otimes(g_jg_j\T)\bigr)+\lambda\Id_{md}.
\end{equation}
При укладке по строкам порядок множителей Кронекера меняется.
Формула \eqref{eq:kron} описывает структуру; строить плотную матрицу
размера $md\times md$ для её применения не требуется.

\subsection{Сертификаты полного условного решения}\label{sec:certificate}
Из $Q=-\mathcal L(B-B_\star)$ следуют
\begin{align}
 \|B-B_\star\|\F&\le\frac{\|Q(B)\|\F}{\lambda},\label{eq:solutionbound}\\
 F(B)-F(B_\star)
 &=\ip{Q(B)}{\mathcal L^{-1}Q(B)}
 \le\frac{\|Q(B)\|\F^2}{\lambda}.\label{eq:gapbound}
\end{align}
В частности, для заданной абсолютной точности $\varepsilon_F$ достаточно
проверить $\|Q\|\F^2/\lambda\le\varepsilon_F$.
Оценка с нормированной невязкой $\|Q\|\F/\max(1,\|C\|\F)$ удобна для
диагностики, но сама по себе не равна относительной ошибке матрицы.

Для решения с ограниченным рангом $Q$ обычно не обращается в ноль.
Поэтому \eqref{eq:gapbound} сравнивает его с \emph{полным} минимумом;
это не критерий достижения оптимума на множестве матриц заданного ранга.
Достижение численного допуска и уменьшение $F$ относительно входящего
базиса также следует проверять отдельно.

При $\lambda=0$ единственность полного решения требует инъективности
$\mathcal A$. Если она отсутствует, нужны явные условия идентифицируемости
и выбор решения, например минимальной нормы. Оценки с делением на $\lambda$
в этом случае неприменимы. Если инъективность установлена, вместо $\lambda$
можно использовать доказанную положительную нижнюю границу спектра
$\mathcal A^*\mathcal A$.

\section{Исходный алгоритм последовательного построения}\label{sec:greedy}

\subsection{Точный шаг вдоль компоненты ранга один}
Пусть текущая матрица равна $B$, а добавка имеет вид $\sigma av\T$,
где $\|a\|_2=\|v\|_2=1$. Обозначим
\begin{align}
 R_B(a,v)&=a\T Q(B)v
 =\sum_jN_j(a\T g_j)e_j\T U_jv+\lambda a\T(P-B)v,\label{eq:R}\\
 D(a,v)&=\ip{av\T}{\mathcal L(av\T)}
 =\sum_jN_j(a\T g_j)^2\|U_jv\|_2^2+\lambda.\label{eq:D}
\end{align}
Тогда
\begin{equation}\label{eq:linesearch}
 F(B+\sigma av\T)-F(B)=D(a,v)\sigma^2-2R_B(a,v)\sigma.
\end{equation}
При $\lambda>0$ знаменатель положителен. Точный оптимальный масштаб и
соответствующее уменьшение функционала равны
\begin{equation}\label{eq:gain}
 \sigma_\star=\frac{R_B(a,v)}{D(a,v)},
 \qquad \gamma(a,v)=\frac{R_B(a,v)^2}{D(a,v)}.
\end{equation}
Знак $\sigma_\star$ можно перенести в $a$ или $v$.
Если $R_B=0$, улучшение вдоль этой пары равно нулю; это ещё не доказывает
отсутствие других полезных пар.

Глобально лучшая добавка ранга один определяется величиной
\begin{equation}\label{eq:Gamma}
 \Gamma_1(Q)=\max_{a\ne0,\ v\ne0}
 \frac{\ip{Q}{av\T}^2}{\ip{av\T}{\mathcal L(av\T)}}.
\end{equation}
Без единичных норм регуляризационный член знаменателя равен
$\lambda\|a\|_2^2\|v\|_2^2$. Замена этого выражения на просто $\lambda$
при произвольных $a,v$ была бы ошибкой.
Величина \eqref{eq:Gamma} --- наилучшее уменьшение после оптимизации
масштаба; вычислять её для общего $\mathcal L$ так же просто, как обычную
спектральную норму, оснований нет.

\subsection{Попеременный поиск направлений}
При фиксированном единичном $v$ положим
\begin{equation}\label{eq:Gv}
 G_v=\sum_jN_j\|U_jv\|_2^2g_jg_j\T+\lambda\Id_m,
 \qquad q_v=Qv.
\end{equation}
Максимизация $(a\T q_v)^2/(a\T G_va)$ даёт
\begin{equation}\label{eq:astep}
 G_v\widetilde a=q_v,\qquad a=\widetilde a/\|\widetilde a\|_2.
\end{equation}
Доказательство --- неравенство Коши--Буняковского после замены
$z=G_v^{1/2}a$. Система имеет размер $m\times m$.

При фиксированном единичном $a$ аналогично
\begin{equation}\label{eq:Ha}
 H_a=\sum_jN_j(a\T g_j)^2U_j\T U_j+\lambda\Id_d,
 \quad q_a=Q\T a,
 \quad H_a\widetilde v=q_a,
 \quad v=\widetilde v/\|\widetilde v\|_2.
\end{equation}
Большая система решается через действие
\[
 H_ax=\sum_jN_j(a\T g_j)^2U_j\T(U_jx)+\lambda x.
\]
Вместо явного формирования $H_a$ можно решать расширенную задачу
наименьших квадратов. Пусть $(A_ax)_j=\sqrt{N_j}(a\T g_j)U_jx$.
Тогда $\widetilde v$ минимизирует
\begin{equation}\label{eq:vls}
 \sum_jN_j\|e_j-(a\T g_j)U_jx\|_2^2
 +\lambda\|x-(P-B)\T a\|_2^2.
\end{equation}
Её нормальное уравнение совпадает с \eqref{eq:Ha}. Регуляризация тянет
к $(P-B)\T a$, а не всегда к нулю.

При точных условных решениях значение отношения в \eqref{eq:Gamma}
не убывает. Это не доказывает нахождение глобально лучшей пары или
глобального минимума задачи \eqref{eq:matrix}. Приближённые решения могут
нарушить эту монотонность. Нулевые $\widetilde a,\widetilde v$ не нормируют;
они требуют остановки данной попытки или другой начальной пары.

\subsection{Сжатие факторов без потери матрицы}
После $k$ добавок матрица представлена произведением $LR\T$, где
$L\in\R^{m\times k}$, $R\in\R^{d\times k}$. Выполним сокращённые
QR-разложения
\[
 L=Q_LR_L,\qquad R=Q_RR_R,
\]
а затем SVD небольшой матрицы $R_LR_R\T$:
\[
 R_LR_R\T=\widehat A\Sigma\widehat V\T.
\]
Получим
\begin{equation}\label{eq:compact}
 B=A\Sigma V\T,
 \qquad A=Q_L\widehat A,\quad V=Q_R\widehat V,
 \quad A\T A=V\T V=\Id_k
\end{equation}
при невырожденных факторах; нулевые компоненты удаляются.
Если сохраняются все ненулевые компоненты, матрица и функционал не меняются.
\emph{Дополнительное усечение} маленьких, но ненулевых сингулярных значений
может увеличить $F$: оно оптимально в норме Фробениуса, а не обязательно
в норме $\mathcal L$. После усечения нужна проверка исходного функционала.

\subsection{Совместная переоценка масштабов}\label{sec:scales}
Для фиксированных ортонормированных $A=[a_1,\ldots,a_k]$,
$V=[v_1,\ldots,v_k]$ рассмотрим
$B=\sum_{s=1}^k\sigma_sa_sv_s\T$. Тогда
\begin{equation}\label{eq:scalequadratic}
 F(B)=C_0+\sigma\T M\sigma-2h\T\sigma,
\end{equation}
где
\begin{align}
 M_{s\ell}&=\sum_jN_j(a_s\T g_j)(a_\ell\T g_j)
        (U_jv_s)\T(U_jv_\ell)+\lambda\delta_{s\ell},\label{eq:M}\\
 h_s&=\sum_jN_j(a_s\T g_j)I_j\T U_jv_s+\lambda a_s\T Pv_s.\label{eq:h}
\end{align}
Система $M\sigma=h$ даёт точный минимум \emph{в выбранной линейной оболочке
компонент}. При $\lambda>0$ матрица $M$ положительно определена.
Не требуется вычислять $M^{-1}$. Полученные $\sigma_s$ могут быть
отрицательными или неупорядоченными: это сначала коэффициенты линейной
комбинации. После переноса знаков и сортировки их модули становятся
сингулярными значениями соответствующей матрицы.

Этот шаг сильнее независимого сохранения жадно выбранных масштабов, но
оптимизация $k$ масштабов не совпадает с оптимизацией всех $mk$ коэффициентов
матрицы на пространстве правых направлений. Последняя рассмотрена
в разделе~\ref{sec:coarse}.

\subsection{Полная последовательность и остановка}
В варианте \eqref{eq:matrix} начинают с $B_0=0$. Для $k=1,\ldots,r$:
\begin{enumerate}
 \item Вычисляют $e_j$ и $Q(B_{k-1})$.
 \item Выбирают начальное направление $v$; возможны случайная попытка
 или ведущая правая сингулярная компонента $Q$.
 \item Повторяют \eqref{eq:astep}, \eqref{eq:Ha} до согласованного допуска
 или до ограничения числа итераций.
 \item Вычисляют точный масштаб \eqref{eq:gain} и проверяют уменьшение $F$.
 \item Добавляют компоненту, выполняют \eqref{eq:compact} и решают
 систему совместных масштабов.
\end{enumerate}
Относительное уменьшение $\gamma/\max(1,F(B_{k-1}))$ задаёт практический
критерий прекращения добавок. Малое улучшение \emph{найденной} пары не
сертифицирует малость \emph{лучшего возможного} улучшения.
Аналогично стабилизация $|v_{t+1}\T v_t|$ показывает стабилизацию итераций,
но не глобальную оптимальность. Полезно отдельно хранить невязку условной
линейной системы, число попеременных шагов, фактическое уменьшение и причину
остановки.

Для малоранговой поправки ту же схему применяют к переменной $\Delta$,
заменяя $I_j$ на $I_j-U_jP\T g_j$, а регуляризационную цель $P$ --- на ноль.
Начальное значение $\Delta_0=0$ соответствует допустимому базису $P$.

\section{Спектральная модель и её границы}\label{sec:isotropic}

\subsection{Точное решение при изотропии}
Предположим, что
\begin{equation}\label{eq:isotropy}
 H_j=\eta_j\Id_d,\qquad \eta_j\ge0,
 \qquad G=\sum_jN_j\eta_jg_jg_j\T+\lambda\Id_m.
\end{equation}
Тогда $\mathcal L(X)=GX$, $G\succ0$ и
\begin{equation}\label{eq:whiten}
 F(B)=C_0-\|W\|\F^2+\|G^{1/2}B-W\|\F^2,
 \qquad W=G^{-1/2}C.
\end{equation}
Умножение слева на обратимую матрицу сохраняет ранг.

\begin{theorem}\label{thm:isotropic}
В условиях \eqref{eq:isotropy} один из глобальных минимумов
\eqref{eq:matrix} имеет вид
\begin{equation}\label{eq:isoopt}
 B_r^{\mathrm{iso}}=G^{-1/2}[G^{-1/2}C]_r,
\end{equation}
где $[W]_r$ --- усечённое SVD с $r$ наибольшими сингулярными значениями.
Если эти значения обозначены $s_1\ge s_2\ge\cdots\ge0$, то
\begin{equation}\label{eq:isotail}
 F(B_r^{\mathrm{iso}})-F(B_\star)=\sum_{k>r}s_k^2,
 \qquad B_\star=G^{-1}C.
\end{equation}
\end{theorem}
\begin{proof}
Замена $Z=G^{1/2}B$ превращает задачу в минимизацию
$\|Z-W\|\F^2$ при $\rank(Z)\le r$. По теореме Эккарта--Янга
минимум достигается на $[W]_r$, а его значение равно квадрату нормы
отброшенного спектрального хвоста~\cite{eckart}. Обратная замена даёт
\eqref{eq:isoopt} и \eqref{eq:isotail}.
При кратных значениях на границе усечения решение может быть неединственным.
\end{proof}

На практике симметричные квадратные корни можно заменить разложением
Холецкого $G=TT\T$ с нижней треугольной $T$:
\begin{equation}\label{eq:choleskywhiten}
 B_r^{\mathrm{iso}}=T^{-\mathsf T}[T^{-1}C]_r.
\end{equation}
Здесь $T^{-1}C$ и умножение на $T^{-\mathsf T}$ вычисляются решением
треугольных систем. Именно такой порядок множителей следует из замены
$Z=T\T B$.

Выбор ранга по
\[
 \frac{\sum_{k>r}s_k^2}{\sum_k s_k^2}\le\varepsilon
\]
при ненулевом знаменателе контролирует долю \emph{неполученного уменьшения
от $B=0$ до полного минимума}. Это не относительная ошибка $F(B_r)/F(B_\star)$
и не оценка размерности истинного EDR-пространства.
Для оптимального $B_r^{\mathrm{iso}}$ имеем
\[
 G^{-1/2}Q(B_r^{\mathrm{iso}})=W-[W]_r,
 \qquad \Gamma_1(Q)=s_{r+1}^2.
\]
Для произвольной текущей $B$ изотропная модель также даёт
$\Gamma_1(Q)=\|G^{-1/2}Q\|_2^2$.

\subsection{Изотропное начальное приближение}
В общем случае естественно выбрать
\begin{equation}\label{eq:eta}
 \eta_j=\frac{\tr(H_j)}d=\frac{\|U_j\|\F^2}d.
\end{equation}
Это минимум $\|H_j-\eta\Id_d\|\F^2$ по скаляру $\eta$.
Матрица \eqref{eq:isoopt} с таким $G$ служит спектральным начальным приближением.
Формула точна для приближённого функционала; её точность для исходного
$F$ требует дополнительных оценок или измерений.

Если $p<d$ и $U_j\ne0$, то $\rank(H_j)\le p<d$.
Следовательно, положительная точная изотропия $H_j=\eta_j\Id_d$ невозможна.
Более того,
\begin{equation}\label{eq:anisofloor}
 \|H_j-\eta_j\Id_d\|_2
 \ge\eta_j\max\{1,d/p-1\}.
\end{equation}
Действительно, у $H_j$ есть нулевое собственное значение, а максимальное
не меньше $\tr(H_j)/p=(d/p)\eta_j$.
Поэтому оптимальность скалярного приближения в норме Фробениуса не означает
его близость в операторной норме при $p\ll d$.

\subsection{Проверяемая оценка возмущения}\label{sec:perturbation}
Пусть $\mathcal M(X)=GX$ и
\begin{equation}\label{eq:rho}
 \rho=\sum_jN_j\|g_j\|_2^2\|H_j-\eta_j\Id_d\|_2,
 \qquad \mu_G=\lambda_{\min}(G),\qquad\theta=\rho/\mu_G.
\end{equation}
Из подчинённых норм следует
$\|(\mathcal L-\mathcal M)X\|\F\le\rho\|X\|\F$.
Если $\theta<1$, то
\begin{equation}\label{eq:equivalence}
 (1-\theta)\ip{X}{\mathcal M X}
 \le\ip{X}{\mathcal L X}
 \le(1+\theta)\ip{X}{\mathcal M X}.
\end{equation}
Использование $\mu_G$ уточняет более консервативное условие
$\rho<\lambda$ из исходников, поскольку $\mu_G\ge\lambda$.

Пусть $Z_0=G^{-1/2}uv\T$, где $(u,v)$ --- ведущая сингулярная пара
$G^{-1/2}Q$. После точной оптимизации масштаба по исходному $F$
это направление удовлетворяет
\begin{equation}\label{eq:proposalbound}
 \frac{\gamma(Z_0)}{\Gamma_1(Q)}\ge\frac{1-\theta}{1+\theta}
 \quad\text{при }\Gamma_1(Q)>0,
 \qquad
 \Gamma_1(Q)\le\frac{\|G^{-1/2}Q\|_2^2}{1-\theta}.
\end{equation}
Для доказательства достаточно заменить знаменатель отношения
\eqref{eq:Gamma} двумя границами \eqref{eq:equivalence} и применить
изотропный максимум. Левая компонента должна быть преобразована
в $G^{-1/2}u$; использование одного лишь $u$ меняет направление.

Эти оценки относятся к одной \emph{свободной добавке} ранга один.
Она может увеличить текущий ранг. Оценки не утверждают, что все замены
компонент при сохранении ранга исчерпаны.
Если $\theta\ge1$, знаменатель верхней оценки использовать нельзя.
Спектральное начальное приближение остаётся допустимой эвристикой.
Оценка $\rho$ может быть грубой из-за отдельного суммирования норм;
проверка истинной относительной операторной ошибки иногда даёт меньшую
величину, но требует собственного надёжного способа сертификации.

\subsection{Стоимость спектрального начального приближения}
Для явных плотных $U_j$ один проход позволяет получить
$c_j=U_j\T I_j$ и $\|U_j\|\F^2$ за $O(Jpd)$ операций.
Далее накопление $C,G$ и разложения требуют
\begin{equation}\label{eq:isocost}
 O(Jpd+Jmd+Jm^2+m^2d+m^3)
\end{equation}
операций. Дополнительная память $O(md+m^2+pd)$ достаточна при последовательной
обработке центров; последний член --- рабочий буфер, если он нужен.
Это оценка \emph{сверх хранения входных данных}. Сам массив $U$ занимает
$Jpd$ чисел. При пакетной обработке буферы увеличиваются с размером пакета.
Переход к случайному приближённому SVD стоит рассматривать только после
измерения затрат: при малом $m$ точное сокращённое SVD матрицы $m\times d$
может быть достаточно дешёвым.

\section{Структурные ограничения и восстановление пространства}\label{sec:limits}

\subsection{Неустранимый порог при усечении самой матрицы}
Поскольку $PP\T=\Id_m$, у $P$ ровно $m$ сингулярных значений, равных единице.
По теореме Эккарта--Янга
\begin{equation}\label{eq:floor}
 \min_{\rank(B)\le r}\|B-P\|\F^2=m-r,
 \qquad F(B)\ge\lambda(m-r).
\end{equation}
Если $F(P)<\lambda(m-r)$, ни один допустимый кандидат \eqref{eq:matrix}
не улучшит функционал относительно $P$.
Это ограничение постановки, а не недостаток конкретного поиска направлений.

Для глобального ограниченного минимума $B_r^\star$ из \eqref{eq:energy}
и \eqref{eq:bounds} следует
\begin{equation}\label{eq:ranktail}
 F(B_r^\star)-F(B_\star)
 \ge\lambda\sum_{k>r}\sigma_k(B_\star)^2.
\end{equation}
Если $\|B_\star-P\|_2\le\delta<1$, то
$\sigma_k(B_\star)\ge1-\delta$ и
\begin{equation}\label{eq:nearprior}
 F(B_r^\star)-F(B_\star)
 \ge\lambda(m-r)(1-\delta)^2.
\end{equation}
Близость полного решения к ортонормированному базису, таким образом,
не поддерживает предположение о быстро убывающем спектре самой $B_\star$.
Спектр поправки $B_\star-P$ может убывать иначе.

\subsection{Недостающие направления не возникают при ортонормировании}
Для любой положительно полуопределённой матрицы $S_g$ имеем
\[
 \rank(B\T S_gB)\le\rank(B)\le r.
\]
При $r<m$ из одной $B_r$ нельзя получить $m$ определённых ею ненулевых
структурных направлений. Произвольные векторы нулевого пространства SVD
не являются восстановленными направлениями EDR.
В действующей реализации недостающие направления дополняются из текущего
$P$ после проекции на дополнение к найденным правым направлениям.
Это вычислительно определённая процедура, но её качество не следует из
минимальности $F(B_r)$.

Для поправки $P+\Delta$ порога \eqref{eq:floor} нет: $\Delta=0$ допустима.
При точной минимизации её функционал не больше $F(P)$.
Однако полный ранг результата не гарантирован. Достаточное условие
$\|\Delta\|_2<1$ даёт
\begin{equation}\label{eq:rankpreserve}
 \sigma_m(P+\Delta)\ge1-\|\Delta\|_2>0.
\end{equation}
Это достаточная проверка, а не обязательное ограничение всякой полезной
поправки.

\subsection{Что меняется при извлечении базиса}
Если $B$ имеет полный ранг строк, можно представить $B=TQ$,
где $Q Q\T=\Id_m$ и $T$ обратима. Тогда
\begin{equation}\label{eq:coeftransform}
 B\T g_j=Q\T(T\T g_j).
\end{equation}
Замена $g_j$ на $T\T g_j$ сохраняет предсказания и ошибку данных.
Точная повторная локальная оценка коэффициентов на $Q$ может только уменьшить
эту ошибку. Численное отсечение малых значений в локальных задачах требует
проверки погрешности отдельно.
Штраф $\|Q-P\|\F^2$ при этом не обязан совпадать с $\|B-P\|\F^2$.
Следовательно, уменьшение исходного $F$ для сырой $B$ не переносится
автоматически на ортонормированный $Q$.

\subsection{Условная связь матричной ошибки с ошибкой проектора}
Пусть полный условный минимум имеет
$s=\sigma_m(B_\star)>0$, а $\widehat B$ удовлетворяет
$\|\widehat B-B_\star\|_2<s$. Тогда $\widehat B$ тоже имеет ранг $m$.
Для проекторов на пространства строк этих матриц
\begin{equation}\label{eq:projectorbound}
 \|\widehat\Pi-\Pi_\star\|_2
 \le\frac{\|\widehat B-B_\star\|_2}{s},
 \qquad
 \|\widehat\Pi-\Pi_\star\|\F
 \le\frac{\sqrt2\|\widehat B-B_\star\|\F}{s}.
\end{equation}
Для вывода заметим, что
$B_\star(\Id-\widehat\Pi)
=(B_\star-\widehat B)(\Id-\widehat\Pi)$, и применим SVD $B_\star$.
Норма разности проекторов равной размерности выражается через синусы
главных углов. Сочетание с \eqref{eq:solutionbound} даёт условный численный
контроль, но при малом $s$ он становится слабым.

Проектор $\Pi_\star$ в \eqref{eq:projectorbound} относится к минимуму
\emph{зафиксированной задачи}, а не к истинному $\Pi_\circ$.
Для статистической теории требуется отдельная цепочка:
\[
 \begin{gathered}
 \text{ошибки локальных статистик}
 \longrightarrow\text{ошибка условного решения}\\
 \longrightarrow\text{ошибка пространства}
 \longrightarrow\text{поведение внешнего цикла}.
 \end{gathered}
\]
Ни точное условное решение, ни малая ошибка данных не заменяют эту цепочку.

\subsection{Что уже делает код и что пока не установлено}
На дату подготовки документа \src{ADP/solver/SVD.py} содержит явные
варианты усечения матрицы и малоранговой поправки. Он работает на CPU,
решает малые положительно регуляризованные задачи по $v$ методом Холецкого
при $d\le128$, а остальные --- через LSMR с проверкой невязки нормального
уравнения. Порог $128$ является настройкой данной реализации, а не
математическим пределом. Усечение матрицы завершается дополнением базиса
из $P$; вариант поправки проверяет ранг и ортонормирует результат.

Сохранённое парное сравнение полного обучения содержит $350$ успешных
запусков: $100$, $50$ и $25$ пар для трёх конфигураций. В большой конфигурации
$(n,d,m,J,p)=(800,50,2,800,10)$ медиана ошибки проектора составила $0{,}118$
для варианта SVD с $r=1$ и $0{,}0464$ для существующего LSMR;
медианы времени --- $24{,}8$ и $37{,}1$ секунды соответственно.
Это сохранённые результаты
\src{experiments/lsmr_vs_svd_fullfit_2026-09-29/REPORT.md},
а не новые измерения. Использовались двойная точность, один поток BLAS
и отдельный процесс на каждый запуск; параметры локализации, допуски
и начальные состояния сохранены рядом с отчётом. Ошибка проектора
определялась как норма Фробениуса его разности с истинным проектором,
делённая на $\sqrt{2m}$. Сравнивались разные внутренние постановки; данные
характеризуют полное обучение и не доказывают превосходства одного
оптимизатора на общем функционале. Успешное завершение запуска также
не равнозначно выполнению критерия сходимости.

Полный условный шаг исследован в отдельном прототипе. Проверка небольшой
линейной задачи подтверждает формулы, но завершённого исследования,
которое обосновывало бы улучшение восстановления на новых данных,
в исходном состоянии проекта нет. Именно поэтому дальнейшие разделы
задают программу развития, а не объявляют её результаты полученными.

\section{Направление I: сочетание SVD со вторым решателем}\label{sec:second}

Цель этого направления --- использовать спектральный этап для поиска
полезного пространства или начального приближения, затем устранить
остаточную ошибку решением исходной условной задачи.
Для задачи с ограничением $\rank(B)\le r$ уточнение всех $md$ неизвестных
снимает ограничение и тем самым меняет постановку.
Для уже выбранной полной задачи спектральная подготовка и уточнение
являются двумя этапами решения одного функционала.

\subsection{Полное уточнение через расширенный оператор}\label{sec:fullsolve}
Для $B=P+\Delta$ положим $b_0=b-\mathcal A P$. Тогда
\begin{equation}\label{eq:fullcorrection}
 \min_\Delta\|\mathcal A\Delta-b_0\|_2^2+\lambda\|\Delta\|\F^2.
\end{equation}
LSMR решает такую задачу через прямое и сопряжённое действия оператора;
явная большая нормальная матрица не нужна~\cite{lsmr}.
Для этой записи параметр регуляризации LSMR равен $\sqrt\lambda$.
Исходные параметры и смысл начального приближения следует сверять с
документацией конкретной реализации~\cite{scipy}.

Если спектральный этап дал произвольную $B_0$, наиболее прозрачное
уточнение --- найти $D$ из задачи
\begin{equation}\label{eq:shifted}
 \min_D
 \left\|
 \begin{bmatrix}\mathcal A\\ \sqrt\lambda\Id\end{bmatrix}D-
 \begin{bmatrix}b-\mathcal A B_0\\\sqrt\lambda(P-B_0)\end{bmatrix}
 \right\|_2^2,
 \qquad B=B_0+D.
\end{equation}
Она точно совпадает с исходным полным функционалом.
Сдвинуть лишь невязку данных и добавить штраф $\lambda\|D\|\F^2$ было бы
неверно: центр регуляризации изменился бы с $P$ на $B_0$.
После уточнения проверяют $Q(B)$, $F(B)\le F(P)$ с поправкой на округление
и полный ранг, если требуется извлечь $m$-мерный базис.

\subsection{Точное решение на найденном пространстве}\label{sec:coarse}
Пусть SVD нашло ортонормированные правые направления
$V\in\R^{d\times q}$, $V\T V=\Id_q$. На линейном пространстве
\[
 \mathcal S_V=\{YV\T:Y\in\R^{m\times q}\}
\]
можно минимизировать исходный $F$ точно по всем $mq$ коэффициентам:
\begin{equation}\label{eq:projected}
 \lambda Y+\sum_jN_j(g_jg_j\T)Y(U_jV)\T(U_jV)=CV.
\end{equation}
Это положительно определённая система размера $mq$.
При малых $m,q$ допустимо сформировать её небольшую плотную матрицу.
Матрицы $U_jV$ можно вычислять пакетами и накапливать систему без хранения
всех $J$ блоков. Полученный $B_V=YV\T$ является наилучшим приближением
$B_\star$ в норме $\mathcal L$ на $\mathcal S_V$.
Если начать с матрицы, уже лежащей на $\mathcal S_V$, решение
\eqref{eq:projected} не увеличит её функционал.

Но при $q<m$ матрица $B_V$ остаётся неполного ранга.
\proposal Для полного шага предпочтительно рассмотреть
\emph{аффинное} пространство
\[
 B=P+YV\T.
\]
Тогда левая часть \eqref{eq:projected} сохраняется, а правая становится
\begin{equation}\label{eq:affine}
 \lambda Y+\sum_jN_j(g_jg_j\T)Y(U_jV)\T(U_jV)=Q(P)V.
\end{equation}
Поскольку $Y=0$ допустима, точный результат имеет $F(B_V)\le F(P)$.
Это малоранговая поправка на выбранном пространстве; последующее полное
уточнение снимает и это ограничение. Сама аффинная запись не гарантирует
ранг, поэтому проверка \eqref{eq:rankpreserve} или прямой контроль
$\sigma_m(B_V)$ всё равно нужны.

\subsection{Уточнение методом сопряжённых градиентов}
В силу леммы~\ref{lem:spd} нормальное уравнение допускает метод сопряжённых
градиентов. После первого этапа вычисляют $R=C-\mathcal L(B_V)$ и решают
\begin{equation}\label{eq:refine}
 \mathcal L(D)=R,\qquad B=B_V+D.
\end{equation}
Дешёвый предобусловливатель задаётся изотропной моделью:
\begin{equation}\label{eq:precondition}
 \mathcal M^{-1}(R)=G^{-1}R.
\end{equation}
Одно разложение $G$ используется во всех итерациях, пока фиксированы
величины, определяющие $G$. При \eqref{eq:equivalence} спектральная
обусловленность \emph{симметрично предобусловленного} оператора удовлетворяет
\begin{equation}\label{eq:condition}
 \kappa\bigl(\mathcal M^{-1/2}\mathcal L\mathcal M^{-1/2}\bigr)
 \le\frac{1+\theta}{1-\theta}.
\end{equation}
Для такого оператора классическая оценка ошибки метода сопряжённых
градиентов после $t$ шагов имеет множитель
$2((\sqrt\kappa-1)/(\sqrt\kappa+1))^t$ в энергетической норме.
Она описывает точную арифметику и фиксированный положительно определённый
предобусловливатель. В конечной точности нужны фактические невязки.

LSMR удобно применять к расширенной задаче; предобусловленный метод
сопряжённых градиентов --- к симметричному положительно определённому
оператору. Оба пути нужно сравнивать по числу действий $U,U\T$, невязкам,
времени и памяти. Работа с нормальным уравнением может усилить проявление
плохой обусловленности; отсутствие явного формирования матрицы само по
себе этого не устраняет.

\subsection{Важное уточнение к ортогональности невязки}
Из \eqref{eq:projected} или \eqref{eq:affine} следует $RV=0$.
Это ортогональность невязки к $\mathcal S_V$ по Фробениусу, но из неё
\emph{не следует} $DV=0$ для решения \eqref{eq:refine}.
Оператор $\mathcal L$ может связывать найденное пространство с его
дополнением. Например, при $m=1$ и двух координатах возьмём
\[
 \mathcal L\leftrightarrow
 \begin{pmatrix}2&1\\1&2\end{pmatrix},
 \quad V=\begin{pmatrix}1\\0\end{pmatrix},
 \quad R=(0,1).
\]
Тогда $RV=0$, но $D=(-1/3,2/3)$ и $DV\ne0$.
Следовательно, просто запретить уточнению изменять найденные координаты
означало бы изменить задачу. Для удаления уже разрешённых компонент
из процесса Крылова нужна согласованная проекция в метрике $\mathcal L$.

\subsection{Расширение и повторное использование пространства}
\proposal До полного уточнения можно несколько раз расширить $V$
правыми направлениями $G^{-1/2}R$, ортогонализовать объединение и снова
решить \eqref{eq:affine}. Вложенность аффинных пространств гарантирует
неувеличение $F$ при точных решениях. Выгодность этой схемы определяется
стоимостью построения направлений и малых систем, а не одной их размерностью.

Между соседними внешними шагами можно сохранять небольшое число полезных
правых направлений. После изменения $U,g,N,\lambda$ требуется заново
считать $U_jV$, проекционную систему и невязку. При повороте базиса $P$
коэффициенты $g_j$ должны быть согласованно преобразованы.
Нельзя переносить старый предобусловливатель как будто оператор не изменился.

\proposal Ввести ограниченный бюджет первого этапа: завершать расширение
по достижении заданной невязки, прекращении полезного уменьшения на один
проход по $U$ или исчерпании памяти; остаток решать полным методом.
Бюджет фиксируется до проверочных запусков.
Если итоговый второй этап сохраняет ту же полную задачу и тот же сертификат,
сравнение остаётся сравнением способов решения одного функционала.

\section{Направление II: пакетная структура и GPU}\label{sec:gpu}

\subsection{Два вида параллелизма}
При фиксированных $B,g$ локальные операции для разных центров $j$
независимы до суммирования их вкладов. Их можно выполнять параллельно.
Напротив, жадные компоненты $k$ используют разные невязки
$Q(B_{k-1})$: следующая зависит от предыдущей.
Независимый запуск всех компонент от одного $Q$ не воспроизводит алгоритм
раздела~\ref{sec:greedy} и может найти повторяющиеся направления.

Для одновременной работы с несколькими направлениями нужна
\emph{совместная блочная подзадача}. Она сохраняет исходный функционал,
но меняет траекторию приближённого поиска. Результаты такого варианта
следует проверять отдельно от простого переноса арифметики на GPU.

\subsection{Совместная добавка ранга не выше q}\label{sec:block}
Пусть текущая матрица --- $B$, $E_j=I_j-U_jB\T g_j$, $Z=P-B$.
Ищем добавку $L_bV\T$, где
$L_b\in\R^{m\times q}$, $V\in\R^{d\times q}$:
\begin{equation}\label{eq:blockobj}
 \min_{L_b,V}\Phi(L_b,V)
 =\sum_jN_j\|E_j-U_jV L_b\T g_j\|_2^2
   +\lambda\|L_bV\T-Z\|\F^2.
\end{equation}
Точно $F(B+L_bV\T)=\Phi(L_b,V)$.
При $q=1$ это поиск одной добавки со свободным масштабом.
После добавки ранг не превышает $\rank(B)+q$.

Пусть сначала $V\T V=\Id_q$, $Y_j=U_jV$, $T_j=Y_j\T Y_j$.
Условная задача по $L_b$ имеет уравнение
\begin{equation}\label{eq:blockleft}
 \sum_jN_jg_jg_j\T L_bT_j+\lambda L_b
 =\sum_jN_jg_jE_j\T Y_j+\lambda ZV.
\end{equation}
При малых $m,q$ это плотная система по $mq$ неизвестным.
Если $V$ не ортонормирована, регуляризационный член слева равен
$\lambda L_b(V\T V)$, а не $\lambda L_b$.

Для фиксированного $L_b$ положим $\alpha_j=L_b\T g_j$ и обозначим
переменную правого фактора $X\in\R^{d\times q}$. Её уравнение:
\begin{equation}\label{eq:blockright}
 \sum_jN_jU_j\T U_jX\alpha_j\alpha_j\T
 +\lambda X(L_b\T L_b)
 =\sum_jN_jU_j\T E_j\alpha_j\T+\lambda Z\T L_b.
\end{equation}
Правый фактор содержит $dq$ неизвестных. Его \emph{столбцы связаны}
матрицами $\alpha_j\alpha_j\T$ и $L_b\T L_b$; это не $q$ независимых задач
по $d$ неизвестных.
При $\lambda>0$ положительная определённость по $X$ обеспечена, если
$L_b$ имеет полный ранг столбцов. Если ранг потерян, этот вывод не действует:
следует удалить зависимые компоненты или явно решать вырожденную задачу.

Можно сохранить постановку наименьших квадратов:
\begin{equation}\label{eq:blockaug}
 \mathcal T_{L_b}(X)=
 \left(\{\sqrt{N_j}U_jX\alpha_j\}_j,
       \sqrt\lambda X L_b\T\right),
 \qquad t=\left(\{\sqrt{N_j}E_j\}_j,\sqrt\lambda Z\T\right).
\end{equation}
Минимизируется $\|\mathcal T_{L_b}(X)-t\|^2$; сопряжённое действие равно
\[
 \mathcal T_{L_b}^*(\{z_j\},W)
 =\sum_j\sqrt{N_j}U_j\T z_j\alpha_j\T+\sqrt\lambda W L_b.
\]
Вместо стандартной регуляризации скалярным $\sqrt\lambda\Id$
здесь действует правое умножение на $L_b\T$.
Замена этого блока на единичный меняет функционал.
LSMR видит один вектор длины $dq$, но действия над ним используют матрицы.
Это отличается от блочного метода Крылова с несколькими правыми частями.

После решения по $X$ можно сделать $X=\widetilde V R$, где
$\widetilde V\T\widetilde V=\Id_q$, и заменить $L_b$ на $L_bR\T$.
Произведение $L_bX\T$ сохранится. При точных условных решениях $\Phi$
не возрастает. Итог проверяют по исходному $F$.

\subsection{Почему матричные действия удобны для GPU}
В последовательном варианте многократно вычисляется $U_jv$.
В блочном варианте тяжёлая операция --- $U_jV$, одновременно для $q$
столбцов. При регулярной укладке центров это реализуется как крупное
умножение плоского массива $(Jp)\times d$ на $d\times q$ либо как пакет
матричных умножений. Библиотека cuBLAS предоставляет соответствующие
пакетные операции~\cite{cublas}.

Умножение одной плотной матрицы на $q$ столбцов требует порядка
$2Jpdq$ арифметических операций. Если $U$ читается из памяти один раз,
а остальные затраты пока пренебрежимо малы, отношение операций к объёму
чтения $U$ в двойной точности составляет примерно $q/4$ операций на байт.
Для одного вектора это примерно $1/4$.
Такой расчёт объясняет возможную пользу повторного использования элементов
$U$, но не предсказывает фактическое ускорение: влияют размеры, временные
буферы, устройство памяти, число запусков и производительность двойной
точности конкретного GPU. При $r=1$ или $2$ основной параллелизм может
идти по центрам, а не по числу направлений.

\subsection{Пакетная реализация с ограниченной памятью}
Для пакета из $s$ центров хранение явных операторов требует $spd$ чисел,
а $U_jV$ --- $spq$. Остальные буферы включают невязки $sp$ и факторы
размеров $dq,mq$. Типичная дополнительная память имеет порядок
\begin{equation}\label{eq:gpumemory}
 O(spd+spq+sp+dq+mq+(mq)^2),
\end{equation}
где $(mq)^2$ нужен лишь при плотном формировании малой системы.
Если $U$ уже постоянно размещён на устройстве, $spd$ учитывается в
хранении входа, а не повторно в каждом пакете.
Размер пакета выбирается по измеренному объёму доступной памяти.
Не следует создавать тензор размера $(J,p,m,d)$ или нормальную матрицу
размера $dq\times dq$ для общего действия \eqref{eq:blockright}.

При одинаковом $p$ удобна регулярная укладка $(s,p,d)$.
При переменной длине локальных блоков можно группировать центры по длине
или применять списки соседей и разреженное представление. При этом строки
$U_j$ в нашей постановке соответствуют пробным направлениям, а не
автоматически числу локальных наблюдений. Переход к исходным соседям
означает изменение представления статистик и требует отдельной сверки.

Вклады центров накапливаются в небольшие матрицы и векторы. Разбиение суммы
на пакеты не меняет её математически, но меняет порядок округления.
Сравнение CPU и GPU должно учитывать эту погрешность и проверять
сопряжённость действий. Не следует копировать весь $U$ между CPU и GPU
на каждом шаге Крылова. Если $U$ не помещается на устройстве, стоимость
повторной передачи пакетов входит в полное время решения.

\subsection{Дополнительные возможности и условия проверки}
\proposal На первом этапе переноса оставить те же направления, допуски
и функционал, распараллелив действия по $j$. Затем отдельно исследовать
совместные блоки направлений. Такой порядок позволяет понять источник
изменения результата.

Если $U_j$ имеет подтверждённое факторизованное или разреженное
представление, следует применять его прямо, а не сначала собирать плотный
тензор. Разреженность локальных весов сама по себе не доказывает
разреженность $U_j$ или нормального оператора.

\proposal Несколько GPU могут распределить центры между устройствами,
после чего суммировать частичные вклады в $Q$ и в малые системы.
В каждой итерации второго решателя потребуется обмен результатами
суммирования. Выгодность определяется соотношением локальной работы и
объёма обмена; при малой задаче дополнительные устройства могут не окупиться.

Сначала проверяется двойная точность. Смешанная точность --- отдельное
приближение: оно требует контроля невязки и ранга в двойной точности,
а при необходимости повторного уточнения. Само наличие GPU не обосновывает
ослабление численных допусков.

\section{Направление III: улучшение самого алгоритма}\label{sec:algorithm}

\subsection{Два уровня выбора начального приближения}
Нужно различать выбор начального внешнего ADP-базиса $P$ и начальных
внутренних факторов при фиксированных статистиках.
Для внешнего базиса локальные линейные оценки градиента и SVD их матрицы
используют доступный сигнал данных. Вес $\sqrt{N_j}$, регуляризация
локальных задач и выбор полосы влияют на статистическую процедуру.
Изменение этих правил требует проверки полного обучения на независимых
данных. Даже исходный базис высокого качества не устраняет потери
направлений при последующем усечении $\rank(B)<m$.

Для внутреннего этапа естественны:
\begin{enumerate}
 \item спектральное начальное приближение \eqref{eq:isoopt} с
 $\eta_j=\|U_j\|\F^2/d$;
 \item ведущие направления преобразованной невязки $G^{-1/2}Q$;
 \item небольшой набор направлений предыдущей условной задачи после
 обновления оператора;
 \item несколько воспроизводимых независимых начальных попыток.
\end{enumerate}
Во втором варианте левый фактор преобразуют вместе с правым, как
в \eqref{eq:proposalbound}. В третьем заново оценивают пригодность старого
пространства. В четвёртом выбирают попытку по реальному уменьшению $F$,
а не только по длине градиента.

\subsection{Выбор направления по уменьшению, а не только по градиенту}
Ведущая пара SVD $Q$ максимизирует $|a\T Qv|$ при единичных нормах.
Но лучшая пара для шага максимизирует отношение \eqref{eq:Gamma}:
направление с большим числителем может иметь ещё большую кривизну.
Это объясняет мотивацию преобразования $G^{-1/2}Q$ и точной оптимизации
масштаба по исходному знаменателю.

\proposal Сформировать несколько кандидатов: ведущие пары $Q$, пары
$G^{-1/2}Q$, направления прошлого шага и небольшое число случайных пар.
Для каждой пары вычислить \eqref{eq:gain}; только лучшие подвергать
дорогому попеременному уточнению. Прямые действия $U_jV$ для группы кандидатов
можно объединить. Такой отбор может сократить число безрезультатных
линейных решений, но его дополнительный проход по $U$ должен учитываться.

Проекция $Q(\Id-VV\T)$ удобна для получения новых линейно независимых
правых направлений. Однако при общей $\mathcal L$ условие $v\perp V$
не является свойством глобально лучшей добавки. Оно ограничивает поиск.
Поэтому её можно применять для начальных кандидатов или расширения
пространства, после чего совместно переоценивать коэффициенты.
Жёстко запрещать уточнению менять существующие направления без анализа
функционала нельзя.

\subsection{Совместная оптимизация всех факторов}
Вместо накопления независимых компонентов можно сразу записать
$B=LR\T$, $L\in\R^{m\times r}$, $R\in\R^{d\times r}$.
При невязках $e_j=I_j-U_jR L\T g_j$ градиенты равны
\begin{align}
 \nabla_L F&=-2\sum_jN_jg_je_j\T U_jR
              +2\lambda\bigl(L(R\T R)-PR\bigr),\label{eq:gradL}\\
 \nabla_R F&=-2\sum_jN_jU_j\T e_j(g_j\T L)
              +2\lambda\bigl(R(L\T L)-P\T L\bigr).\label{eq:gradR}
\end{align}
При фиксации одного фактора задача по другому квадратична.
Точное попеременное решение блоков не увеличивает $F$, но совместная
задача невыпукла. Уменьшение функционала не доказывает глобальную сходимость
или сходимость самих факторов.

Неопределённость масштаба $L\mapsto LS$, $R\mapsto RS^{-\mathsf T}$
при обратимой $S$ можно уменьшить уравновешиванием:
\[
 LR\T=A\Sigma V\T,
 \qquad L\leftarrow A\Sigma^{1/2},\quad R\leftarrow V\Sigma^{1/2}.
\]
Это сохраняет матрицу и функционал. Нельзя заменять исходный штраф
$\lambda\|LR\T-P\|\F^2$ суммой штрафов факторов без признания изменения задачи.

Римановы методы на многообразии матриц \emph{фиксированного} ранга также
могут использовать \eqref{eq:gradL}--\eqref{eq:gradR} и матричные действия.
Но множество $\rank(B)=r$ не включает границу $\rank(B)<r$.
Для решения задачи с ограничением $\rank(B)\le r$ нужны отдельные правила
изменения ранга и обработки вырождения. Одно указание многообразия
не превращает невыпуклую задачу в глобально решаемую.

\subsection{Выбор ранга и критерии остановки}
Размерность модели $m$ и вычислительный ранг $r$ имеют разные смыслы.
Спектральный хвост изотропной модели может предлагать $r$, но при общей
структуре $H_j$ его ошибка по исходному $F$ должна проверяться.
Малое отдельное $s_{r+1}$ не контролирует сумму всех отброшенных $s_k^2$.
Если знаменатель энергетического отношения равен нулю, это отдельный
вырожденный случай.

Для полного уточнения используется \eqref{eq:gapbound}.
Для свободного расширения на одну компоненту при $\theta<1$ используется
верхняя оценка \eqref{eq:proposalbound}. При отсутствии такой оценки малый
найденный выигрыш является лишь диагностикой данного поиска.
Для фиксированного ранга дополнительно проверяется первая производная
в допустимых касательных направлениях; полная невязка $Q$ не обязана
обращаться в ноль.

\proposal В варианте, сохраняющем усечение самой $B$, разрешать его как
окончательный шаг только при подтверждённой приемлемости спектрального
хвоста и функционала. Если $F(P)<\lambda(m-r)$, сразу переходить к
поправке или полному шагу. Такой переход --- явно новая составная процедура,
а не скрытая оптимизация старой постановки.

\subsection{Допуск линейного решения из оценки реального шага}
\proposal Вместо ослабления допуска только по изменению направления
использовать достаточное условие уменьшения функционала.
Для приближённой поправки $D$ в полной задаче определим
\[
 s=Q(B)-\mathcal L(D).
\]
Точное тождество даёт
\begin{equation}\label{eq:inexactdescent}
 F(B+D)-F(B)
 =-\ip{D}{\mathcal L(D)}-2\ip{s}{D}
 \le-\lambda\|D\|\F^2+2\|s\|\F\|D\|\F.
\end{equation}
Следовательно, если $D\ne0$ и
$\|s\|\F\le\eta\lambda\|D\|\F/2$ при $0\le\eta<1$, то
\begin{equation}\label{eq:safedescent}
 F(B+D)-F(B)\le-(1-\eta)\lambda\|D\|\F^2.
\end{equation}
Это достаточная, возможно консервативная проверка.
Тот же вывод применим к условной квадратичной задаче с её собственным
оператором и нижней границей спектра. Он не переносится автоматически
на последовательность нормированных факторов и последующее усечение.
Итоговый сертификат точности сохраняется отдельно; гарантированный спуск
не означает уже малой ошибки решения.

\subsection{Более точные предобусловливатели и границы их смысла}
Изотропный $G$ может быть слишком грубым при $p\ll d$.
\proposal Рассмотреть диагональ векторизованного оператора:
\begin{equation}\label{eq:diagprecond}
 D_{as}=\lambda+\sum_jN_jg_{ja}^2\sum_{\ell=1}^p U_{j\ell s}^2,
 \quad a=1,\ldots,m,\quad s=1,\ldots,d.
\end{equation}
Она положительна при $\lambda>0$ и требует $md$ чисел.
Другой вариант --- $d$ малых блоков размера $m\times m$:
\[
 G_s=\lambda\Id_m+\sum_jN_j\|U_j(:,s)\|_2^2g_jg_j\T.
\]
Их хранение стоит $O(dm^2)$; сборка и разложения также оплачиваются.
Такие приближения описывают разные масштабы координат, но не учитывают
все межкоординатные связи. Их преимущество перед $G$ требует измерений.

Использование приближённого оператора как предобусловливателя при точных
действиях $\mathcal L$ сохраняет целевой функционал. Использование той же
аппроксимации \emph{вместо} $\mathcal L$ в окончательном решении меняет
задачу или точность её решения. Это различие нужно сохранять во всех вариантах.

\subsection{Другие идеи из материалов}
Выпуклая задача
\[
 \min_B F(B)+\tau\|B\|_*,\qquad\tau\ge0,
\]
где ядерная норма --- сумма сингулярных значений, является самостоятельной
регуляризованной постановкой. Она не является точной заменой ограничения
$\rank(B)\le r$. Её можно исследовать как отдельный метод сравнения.

Случайное построение спектрального пространства~\cite{halko} может быть
полезно при дорогих действиях матрицы или большом размере спектрального
этапа. При малом $m$ вводить его только ради предполагаемого ускорения
не обосновано; нужна оценка числа проходов и ошибки.

Для теории невыпуклой оптимизации перспективно изучить ограничения
кривизны на разностях малоранговых матриц:
\[
 \alpha_r\|X\|\F^2\le\ip{X}{\mathcal L(X)}
 \le\beta_r\|X\|\F^2,\qquad\rank(X)\le2r.
\]
Лемма~\ref{lem:spd} даёт грубые глобальные границы, но для сильных
утверждений о геометрии задачи могут потребоваться существенно более
точные отношения $\beta_r/\alpha_r$ и условия выбора начальных факторов.
Сложность общих взвешенных малоранговых задач не доказывает автоматически
такую же сложность именно оператора \eqref{eq:L}; отдельное утверждение
о его вычислительной сложности здесь не заявляется.

\section{Программа проверки трёх направлений}\label{sec:verification}

\subsection{Что считать сохранением или изменением метода}
\begin{center}
\begin{tabularx}{\linewidth}{@{}>{\raggedright\arraybackslash}p{0.34\linewidth}
 >{\raggedright\arraybackslash}X@{}}
\toprule
Изменение & Математический смысл\\
\midrule
Пакетирование центров, иной порядок умножений, сокращённое QR/SVD
без усечения & Точное преобразование с изменением ошибок округления\\
Предобусловливание, решение треугольных систем вместо обращения
матриц & Численная реорганизация при сохранении исходного оператора\\
Новые начальные факторы, совместные блоки, иной допуск остановки
 & Изменение приближённого поиска; сравнивать при одном функционале\\
Снятие ограничения ранга, переход с ранга матрицы к рангу поправки,
ядерный штраф & Новая постановка и, при включении в ADP, новая процедура\\
Новая локализация, закон пробных направлений, веса начального приближения
 & Изменение статистической процедуры; требуется полное исследование\\
\bottomrule
\end{tabularx}
\end{center}
Внутренние направления $a,v,V$ являются переменными оптимизации.
Пробные направления, формирующие $U_j,I_j$, являются частью построения
статистик. Изменять вторые под видом выбора первых нельзя.

\subsection{Первый уровень: небольшие проверяемые задачи}
До измерения большого обучения необходимо проверить:
\begin{enumerate}
 \item прямое и сопряжённое действия \eqref{eq:A}--\eqref{eq:adjoint},
 операторную форму и градиент против явных малых матриц;
 \item масштаб и уменьшение \eqref{eq:gain}, условные системы и
 совместную переоценку масштабов;
 \item изотропную формулу, спектральный хвост и правильный порядок
 треугольных преобразований;
 \item проекционные системы \eqref{eq:projected}, \eqref{eq:affine},
 блочную систему и сопряжённость \eqref{eq:blockaug};
 \item градиенты факторов, сертификаты ошибки, сохранение произведения
 при QR и поведение при потере ранга;
 \item согласованность пакетов и CPU/GPU в двойной точности.
\end{enumerate}
Обязательны вырожденные $U_j$, почти зависимые коэффициенты, сильно
неравные массы, малые и большие $\lambda$, отсутствие полезного нового
направления и исчерпание числа итераций. Нулевая масса проверяется
на уровне математического оператора; публичный вход может её отклонять.

\subsection{Второй уровень: одинаковые зафиксированные статистики}
Для направления I сравниваются полный LSMR, спектральное начальное
приближение с полным LSMR и аффинное проекционное решение с полным
предобусловленным уточнением. При общем функционале и одинаковой конечной
точности их решения должны согласовываться с эталоном.
Расходы подготовки спектрального пространства входят во время метода.

Для направления II сначала сравниваются одинаковые действия CPU/GPU,
затем последовательные и совместные блоки при одном целевом функционале.
Для направления III по одному меняются начальное приближение, правило отбора
кандидатов и предобусловливатель. Совместное изменение всех правил
не позволяет установить причину результата.

Для каждого варианта записываются размерности $(J,p,d,m,r,q)$,
массы и $\lambda$, тип чисел, случайные начальные состояния, допуски,
число итераций, число действий $U,U\T$, значение $F$, $\|Q\|\F$,
ошибка относительно эталона, время и максимальная память.
Для GPU дополнительно записываются устройство, размер пакета, стоимость
переноса данных и синхронизации. Учитывается время всего решения,
включая подготовку, а не только одного матричного умножения.
Результаты с ошибками и незавершёнными решениями сохраняются.

\subsection{Третий уровень: полное обучение и новые данные}
Если внутренний уровень прошёл проверку, исследуется полный ADP:
\begin{enumerate}
 \item фиксируются данные и начальные условия в каждой паре сравнения;
 \item набор для выбора метода отделяется от проверочного набора;
 \item параметры метода фиксируются до использования проверочного набора;
 \item ошибка EDR-проектора, ошибки предсказания, внешний критерий
 остановки и внутренние невязки записываются отдельно;
 \item сравниваются время, память, численные отказы и худшие случаи,
 а не только средний удачный запуск.
\end{enumerate}
Для ортонормированных оценок естественна ошибка
$\|\widehat P\T\widehat P-P_\circ\T P_\circ\|\F$
или её заранее выбранная нормированная версия.
Определение метрики сохраняется между вариантами.
Результаты на известных генераторах не являются доказательством
статистической состоятельности для всех многоиндексных моделей.

\subsection{Последовательность развития}
Практический порядок объединяет три направления:
\begin{enumerate}
 \item Зафиксировать исходный функционал и эталон полного условного шага;
 определить, сохраняется ли ограничение ранга окончательного решения.
 \item Проверить аффинное спектральное пространство и полное уточнение.
 Это непосредственно исследует потерю направлений при усечении матрицы.
 \item Реализовать совместные действия по центрам и по нескольким
 кандидатам, затем проверить их на GPU при тех же численных условиях.
 \item Сравнить спектральное начальное приближение, отбор по реальному уменьшению,
 совместную переоценку и варианты предобусловливания.
 \item После проверки условных задач провести отделённое исследование
 полного обучения на новых данных.
\end{enumerate}
Первое направление устраняет ограничение допустимого множества только
при явном переходе к полной постановке. Второе использует структуру
вычислений; третье сокращает неэффективный поиск и улучшает численную
организацию. Их совместная схема имеет вид
\[
 \boxed{\begin{gathered}
 \text{спектральные направления}
 \ \longrightarrow\ \text{совместная малая задача}\\
 \longrightarrow\ \text{полное уточнение}
 \end{gathered}}.
\]
Она является обоснованной программой исследования. Ускорение и улучшение
восстановления должны подтверждаться измерениями на её полной реализации.

\appendix
\section{Происхождение материала и исправления}\label{sec:sources}

Документ объединяет все непустые источники в папке \src{SVD/}.
Повторяющиеся выводы сведены к единой системе обозначений.
\begin{longtable}{@{}>{\raggedright\arraybackslash}p{0.28\linewidth}
 >{\raggedright\arraybackslash}p{0.66\linewidth}@{}}
\toprule
Источник & Содержание в этом документе\\
\midrule
\endhead
\src{SVD/SVD.tex} & Постановка, компонентный алгоритм, точный масштаб,
условные системы, сокращённое QR/SVD и переоценка масштабов:
разделы~\ref{sec:setup}--\ref{sec:greedy}.\\
\src{SVD/SVD_form.tex} & Операторная форма, невязка, изотропное решение,
спектральные критерии, возмущение, совместные факторы и стоимость:
разделы~\ref{sec:operator}, \ref{sec:isotropic}, \ref{sec:algorithm}.\\
\src{SVD/problem.md} & Порог регуляризационного штрафа, потеря размерности,
неприменимость положительной изотропии при $p<d$, ограниченность
условной теории и статистический разрыв: раздел~\ref{sec:limits}.\\
\src{SVD/chat_1.md} & Второй решатель, проекционное решение,
предобусловливание и повторное использование пространства:
раздел~\ref{sec:second}.\\
\src{SVD/chat_2.md} & Пустой файл на дату подготовки; математического
содержания не содержит.\\
\src{SVD/chat_3.md} & Зависимость жадных компонент, совместная добавка,
матричные действия, GPU и оптимизация факторов:
разделы~\ref{sec:gpu}, \ref{sec:algorithm}.\\
\bottomrule
\end{longtable}

Существенные уточнения исходных рассуждений:
\begin{enumerate}
 \item Условные минимумы и монотонность не названы глобальной сходимостью.
 \item Уточнение полным решателем явно отделено от задачи с ограниченным
 рангом; регуляризация при сдвиге начальной точки сохранена.
 \item Из $RV=0$ не сделан неверный вывод о принадлежности поправки
 ортогональному дополнению.
 \item Сохранение матрицы при сжатии отделено от дополнительного усечения,
 которое может увеличить исходный функционал.
 \item Данные $p$ строк локального оператора отделены от числа наблюдений;
 одна векторизованная задача по $dq$ неизвестным отделена от нескольких
 независимых задач Крылова.
 \item Численный ранг и спектральная энергия отделены от размерности EDR
 и относительной ошибки конечного функционала.
 \item Ортонормирование сохраняет предсказания после преобразования
 коэффициентов, но не обязано сохранять регуляризационный штраф.
\end{enumerate}

Аффинная проекционная схема \eqref{eq:affine}, уточнение порога через
$\mu_G$, проверка спуска \eqref{eq:safedescent}, отбор кандидатов по
уменьшению и координатные предобусловливатели изложены как дальнейшие
разработки исходных идей. Их алгебраические свойства выведены здесь;
вычислительное и статистическое преимущество пока является гипотезой.

В исходных диалогах имеются ссылки вида
\src{:chatgpt-content-reference}, не содержащие адресов публикаций.
Они не включены как научные ссылки. Упоминавшиеся там частные теоремы
о других статистических методах и неподтверждённые утверждения об их
скорости не перенесены на рассматриваемый решатель.
Список ниже содержит проверяемые источники используемых общих методов;
формулы конкретной ADP-задачи выведены непосредственно в тексте.

\begin{thebibliography}{9}
\bibitem{eckart}
Эккарт К., Янг Г. Приближение матрицы другой матрицей меньшего ранга.
\emph{Psychometrika}, 1936, т.~1, с.~211--218.
\url{https://doi.org/10.1007/BF02288367}.

\bibitem{lsmr}
Фонг Д., Сондерс М. LSMR: итерационный алгоритм для разреженных задач
наименьших квадратов. \emph{SIAM Journal on Scientific Computing},
2011, т.~33, №~5, с.~2950--2971.
\url{https://arxiv.org/abs/1006.0758}.

\bibitem{scipy}
SciPy. Документация функции \src{scipy.sparse.linalg.lsmr}:
расширенная регуляризованная задача и параметры остановки.
\url{https://docs.scipy.org/doc/scipy/reference/generated/scipy.sparse.linalg.lsmr.html}.
Дата обращения: 29 сентября 2026 года.

\bibitem{halko}
Халко Н., Мартинссон П.-Г., Тропп Дж.
Поиск структуры с помощью случайности: вероятностные алгоритмы
построения приближённых матричных разложений.
\emph{SIAM Review}, 2011, т.~53, №~2, с.~217--288.
\url{https://arxiv.org/abs/0909.4061}.

\bibitem{cublas}
NVIDIA. Документация cuBLAS: пакетное матричное умножение,
раздел \src{cublas<t>gemmStridedBatched}.
\url{https://docs.nvidia.com/cuda/cublas/index.html#cublas-t-gemmstridedbatched}.
Дата обращения: 29 сентября 2026 года.
\end{thebibliography}
\end{document}
